% answers_short.tex —— 1988 年·广东卷·理工农医类 · 简短答案速查
% 由 tools/make_short_answers.py 自 answers.tex 的 \gksolhead/\gkanswer 提取，
% 并与 1988/1988_数学真题_广东理_答案要点.md 逐题比对后生成，勿手改（改源 MD 后重跑）。

\begin{gknotice}
\textbf{使用说明}：本节只给最终答案便于快速核对，完整推导见解析版。答案为编写者独立推导，\textbf{非官方标准答案}。
\end{gknotice}

\gkpart{大题}
\begin{itemize}[leftmargin=2.6em]
  \item[\textbf{1.}] (1) C；(2) B；(3) D；(4) C；(5) B；(6) C；(7) D；(8) C；(9) A；(10) B；(11) A；(12) B；(13) D；(14) B；(15) C；(16) C；(17) D；(18) C；(19) B；(20) B；(21) A；(22) D；(23) A（固定源题面所称右图缺失，字母无法独立核验）；(24) B；(25) D
  \item[\textbf{2.}] (1) $x^2+y^2-4y=0$；(2) $\pi$；(3) $2\sqrt2\left(\cos\dfrac{5\pi}{6}+i\sin\dfrac{5\pi}{6}\right)$；(4) $-1$；(5) $-162$
  \item[\textbf{3.}] (1) $2\,\mathrm{cm}^3$；(2) $PA\perp CD$
  \item[\textbf{4.}] (1) 焦点为 $F_1(-5,0)$、$F_2(5,0)$，离心率 $e=\dfrac53$，渐近线为 $y=\pm\dfrac43x$；(2) $\angle F_1PF_2=90^\circ$
  \item[\textbf{5.}] (1) 当 $0<\theta<\pi$ 时，$\arg u=\dfrac\theta2$；当 $\pi<\theta<2\pi$ 时，$\arg u=\pi+\dfrac\theta2$；(2) 不可能是正实数
  \item[\textbf{6.}] (1) $f^{-1}(x)=\dfrac{a^x+a^{-x}}2$；当 $a>1$ 时反函数定义域为 $[0,+\infty)$，当 $0<a<1$ 时定义域为 $(-\infty,0]$；(2) $S_n<2^n-\left(\dfrac{\sqrt2}{2}\right)^n$
\end{itemize}
